% !TeX root = ../tfg.tex
% !TeX encoding = utf8
%
%*******************************************************
% Capítulo 7. Explicabilidad
%*******************************************************

\chapter{Explicabilidad}\label{cap:explicabilidad}

\section{Layer-wise Relevance Propagation}\label{sec:lrp}

% Puntos que debería cubrir:
% - Descomposición de Taylor profunda: aproximación de primer orden de cada neurona alrededor de
%   un punto raíz.
% - Propiedad de conservación: \sum_i R_i^{(l)} = \sum_j R_j^{(l+1)} = f(x).
% - Reglas LRP-\varepsilon y LRP-z^+ y su derivación.
\begin{definicion}[Relevancia y conservación~\cite{bach2015}]\label{def:lrp-conservacion}
\pendiente{definición}
\end{definicion}

\begin{proposicion}[Regla LRP-$z^+$ como descomposición de Taylor profunda~\cite{montavon2017}]\label{prop:lrp-zplus}
\pendiente{enunciado}
\end{proposicion}
\begin{demostracion}
\pendiente{demostración}
\end{demostracion}

\begin{definicion}[Regla LRP-$\varepsilon$]\label{def:lrp-eps}
\pendiente{definición}
\end{definicion}

\section{Concept Relevance Propagation}\label{sec:crp}

% Puntos que debería cubrir:
% - CRP (Achtibat et al., 2023): condicionar la propagación de relevancia a un concepto (canal o
%   neurona) para obtener mapas de relevancia por concepto.
% - Selección de ejemplos de referencia por relevancia (explicaciones glocales).
% - Uso en STOOD-X para seleccionar las características relevantes (no aplicable directamente al
%   Swin; ver capítulo 11).
\pendiente{CRP}

\endinput
